% ---
% id: receiver/lem:exactbasic
% kind: lem
% title: "Range of the exact block bound"
% summary: >-
%   The exact block bound lies between the analytic bound and one on every
%   photon-number block, and it is non-negative when the multiclick element is
%   positive. That positivity holds at every setting of a setting family but
%   does not follow from the multiclick bound's hypotheses.
% prominence: standard
% status: accepted
% uses: [thm:multiclick-bound, def:exactbnd, lem:fock-monotone, lem:fock-block-spectrum, lem:fock-functor]
% mentions: [def:setting-family, rem:exactpositivity]
% assumes: []
% refs: []
% provenance: newly-derived
% lean:
% ---

The three bounds of the lemma need different hypotheses. The comparison with
the analytic bound and the upper bound hold under the hypotheses of
Theorem~\ref{thm:multiclick-bound} alone. Non-negativity needs the
\tref{multiclick element} itself to be positive, \eqref{eq:mcpos}, which holds at every setting
of a \tref{setting family} (Definition~\ref{def:setting-family}\,(a)), because
there $M_{\mc}$ is a sum of elements of a POVM; those are the only instances
to which this module applies Theorem~\ref{thm:multiclick-bound}. It does not follow
from \eqref{eq:bbkhyp} and \eqref{eq:prefactorform}, and
Remark~\ref{rem:exactpositivity} is the instance that shows so.

\begin{lemma}
\label{lem:exactbasic}
Let the hypotheses of Theorem~\ref{thm:multiclick-bound} hold. Then, for every
$n\ge0$,
\begin{equation}\label{eq:exactbasic}
  \Bnd(n)\ \le\ \Bnd^{\ast}(n)\ \le\ 1 .
\end{equation}
If in addition
\begin{equation}\label{eq:mcpos}
  M_{\mc}\ \succeq\ 0 ,
\end{equation}
then $\Bnd^{\ast}(n)\ge0$ as well, so that $\Bnd^{\ast}(n)\in[0,1]$.
\end{lemma}

\begin{proof}
\emph{Lower comparison.} $\Bnd(n)\le\Bnd^{\ast}(n)$ is \eqref{eq:mcbound}.

\emph{Upper bound.} For $\Bnd^{\ast}(n)\le1$ it suffices that
$M_{\mc}\preceq\id$, which is the operator form of the argument that closes
the proof of Theorem~\ref{thm:multiclick-bound}. From $B_k\preceq B$,
Lemma~\ref{lem:fock-monotone} gives $\Fock(\id-B_k)\succeq\Fock(\id-B)$;
multiplying by $\beta_k>0$, summing, and then using $\beta_k\ge\gamma$
together with $\Fock(\id-B)\succeq0$ (Lemma~\ref{lem:fock-block-spectrum}),
\begin{equation}\label{eq:mcupper}
  \sum_{k=1}^{K}\beta_k\Fock\big(\id-B_k\big)
  \ \succeq\ \Big(\sum_{k=1}^{K}\beta_k\Big)\Fock(\id-B)
  \ \succeq\ K\gamma\,\Fock(\id-B)
  \ \succeq\ (K-1)\gamma\,\Fock(\id-B) ,
\end{equation}
so \eqref{eq:prefactorform} gives $M_{\mc}=\id+(K-1)\gamma\Fock(\id-B)
-\sum_k\beta_k\Fock(\id-B_k)\preceq\id$. Every summand is photon-number block
diagonal (Lemma~\ref{lem:fock-functor}\,(i)), so the $n$-photon block of
$M_{\mc}$ is its compression by $\Pi_n$ and inherits the bound:
$\blk{M}{\eqn}_{\mc}\preceq\id^{(=n)}$, whence $\Bnd^{\ast}(n)\le1$.

\emph{Non-negativity.} Under \eqref{eq:mcpos} the same compression inherits
positivity, since
$\bra{\psi}\Pi_nM_{\mc}\Pi_n\ket{\psi}=\bra{\Pi_n\psi}M_{\mc}\ket{\Pi_n\psi}
\ge0$ for every $\ket\psi$; so $\blk{M}{\eqn}_{\mc}\succeq0$ and
$\Bnd^{\ast}(n)\ge0$.
\end{proof}
